% !PRINTED 142
\weberpage{142}
%<<CONT

Consequently from (6) one obtains the following system of formulae due to \gesperrt{Newton}\footnote[1]{\gesperrt{Newton}, Arithmetica universalis, edit. 's Gravesande, p.~592.}:
\begin{equation}
\begin{aligned}
0&=s_1+a_1,\\
0&=s_2+a_1s_1+2a_2,\\
0&=s_3+a_1s_2+a_2s_1+3a_3,\\
&\ \vdots\\
0&=s_{n-1}+a_1s_{n-2}+a_2s_{n-3}+\cdots+(n-1)a_{n-1}.
\end{aligned}
\tag{7}
\end{equation}

This system of formulae can, however, be continued still further; for, since
\[
 S[f(\alpha)]=0,\qquad S[\alpha f(\alpha)]=0,\qquad S[\alpha^2f(\alpha)]=0,\ldots,
\]
it follows that
\begin{equation}
\begin{aligned}
0&=s_n+a_1s_{n-1}+a_2s_{n-2}+\cdots+n a_n,\\
0&=s_{n+1}+a_1s_n+a_2s_{n-1}+\cdots+a_ns_1,\\
0&=s_{n+2}+a_1s_{n+1}+a_2s_n+\cdots+a_ns_2,\quad\ldots
\end{aligned}
\tag{8}
\end{equation}

By the formulae (7), (8) the problem is now solved of representing, one after another, the functions $s_1,s_2,s_3,\ldots$ up to any height as whole rational functions of the symmetric fundamental functions; for example
\begin{equation}
\begin{aligned}
 s_1&=-a_1,\\
 s_2&=+a_1^2-2a_2,\\
 s_3&=-a_1^3+3a_1a_2-3a_3,\\
 s_4&=+a_1^4-4a_1^2a_2+4a_1a_3+2a_2^2-4a_4,\\
 &\hspace{1.5em}\vdots
\end{aligned}
\tag{9}
\end{equation}
The manner of formation of these expressions shows that the coefficients in these representations are \gesperrt{whole numbers}.

Conversely, by means of the formulae (7) and (8), one can express the symmetric fundamental functions $a_1,a_2,a_3,\ldots$ rationally by the power sums $s_1,s_2,s_3,\ldots$. This representation is, however, far less simple in so far as the coefficients are not whole but fractional numbers, for example
\begin{equation}
\begin{aligned}
 a_1&=-s_1,\\
 2a_2&=+s_1^2-s_2,\\
 6a_3&=-s_1^3+3s_1s_2-2s_3,\\
 &\hspace{1.5em}\vdots
\end{aligned}
\tag{10}
\end{equation}
